\subsubsection{Апаратна оптимізація: реалізація матричного методу засобами OpenCV}

Хоча базова матрична реалізація з використанням контейнерів \texttt{std::vector} та пулу потоків дозволила суттєво підвищити швидкодію, ручне управління операціями множення не дозволяє повною мірою використати апаратні можливості сучасних процесорів (наприклад, векторні SIMD-інструкції AVX2/AVX-512). Для досягнення максимальної продуктивності обчислень було прийнято рішення інтегрувати розроблений математичний апарат із бібліотекою комп'ютерного зору OpenCV (\texttt{processOpenCVBarycentricZoom}).

Головна перевага використання OpenCV полягає в наявності оптимізованого бекенду для операцій лінійної алгебри. Під капотом структури \texttt{cv::Mat} та її перевантажених операторів множення працюють високоефективні виклики до спеціалізованих математичних бібліотек (таких як BLAS, LAPACK або Eigen), які автоматично забезпечують апаратну векторизацію та оптимальне управління кешем на найнижчому рівні.

Алгоритмічний конвеєр OpenCV-реалізації складається з наступних кроків:
\begin{enumerate}
	\item \textbf{Конвертація типів:} Вихідний об'єкт \texttt{QImage} перетворюється на багатоканальну матрицю \texttt{cv::Mat}. Далі, за допомогою функції \texttt{cv::split}, зображення розшаровується на три окремі одноканальні матриці (канали R, G, B), що відповідають плоским масивам $F^{(c)}$ з попереднього методу.
	\item \textbf{Формування матриць ваг:} Нормалізовані вагові коефіцієнти Чебишева для обох осей генеруються за тими ж формулами \eqref{eq:matrix_Wy}, проте одразу записуються у структури \texttt{cv::Mat} типу \texttt{CV\_32F} (32-бітні числа з плаваючою комою). Це дозволяє уникнути накладних витрат на подальше копіювання пам'яті.
	\item \textbf{Оптимізоване множення:} Обчислення результуючої матриці кольору $C^{(c)}$ виконується за допомогою стандартного оператора множення матриць бібліотеки OpenCV, що математично еквівалентно виклику узагальненого множення матриць (GEMM - General Matrix Multiply):
	\begin{equation}\label{eq:opencv_gemm}
		C^{(c)} = W_y \times F^{(c)} \times W_x^T.
	\end{equation}
	Оскільки OpenCV автоматично розпаралелює ці операції (за наявності бекенду OpenMP або TBB), відпадає необхідність у ручному створенні пулу потоків \texttt{std::thread}.
	\item \textbf{Зворотне злиття:} Після незалежного обчислення нових розмірів для кожного колірного каналу, матриці об'єднуються назад у єдину багатоканальну структуру за допомогою функції \texttt{cv::merge}, нормалізуються і конвертуються назад у формат \texttt{QImage} для відображення у графічному інтерфейсі.
\end{enumerate}

Переведення математичної моделі на рейки індустріального стандарту OpenCV не змінило теоретичної асимптотичної складності алгоритму, яка залишилася на рівні $\mathcal{O}(H_2 \cdot W \cdot (H + W_2))$. Однак, завдяки глибокій оптимізації доступу до пам'яті та використанню векторних інструкцій процесора, практичний час виконання (wall-clock time) зменшився. Ця імплементація продемонструвала найкращий баланс між якістю зображення, притаманною барицентричній інтерполяції, та швидкістю обробки, наблизившись до показників вбудованих методів масштабування самої бібліотеки OpenCV.